\begin{abstract}
A UMI-style device carries a wrist camera for the policy and a pose tracker for the labels: SLAM on that camera, a visual-inertial tracker, a VR controller or a fixed external camera. We replace both with one egocentric camera worn on the collector's head. It sees the fiducial ball and the jaw tags on an otherwise electronics-free glove and the printed tiles taped around the workspace, so every frame carries the gripper's pose in a tile-fixed world frame and the jaw width, with no cable, no IMU and no SLAM. At deployment the same camera rides a cheap servo-driven printed arm beside the robot and keeps the gripper in view the way the head kept the hand in view, driven by a linear viewpoint model fitted on one glove session. The labels are checked against the robot's encoders: \SI{6.5}{mm} median tool-point error, \SI{1.6}{\degree} rotation, \SI{93}{}--\SI{100}{\percent} of frames labelled per session against \SI{78}{}--\SI{92}{\percent} for a fixed external camera that the collector's body hides, and a policy trained on head-camera labels instead of encoders pays \SI{3}{\percent} in validation loss. One person records 2.0--3.1 demonstrations per minute against 1.6 by teleoperation, without occupying the robot, so collectors can work in parallel and a new workspace needs only printed tiles.Adding 194 glove segments to 172 teleoperated episodes lowers ACT's loss on held-out teleoperated episodes from 0.318 to 0.263; added to 48 teleoperated episodes it reaches 0.295, the level of 96--172 teleoperated episodes alone, and on the robot the 48-episode co-trained policy succeeds in \grid{X/25} trials against \grid{X/25} for 48 teleoperated episodes alone and \grid{X/25} for the full 172. Against the conventional wrist-camera recipe and a fixed-camera recipe trained from the same demonstrations the head-camera policy \grid{succeeds in X/25 vs X/25 and X/25}. The glove costs \$22, the camera \$122 and the arm about \$120 in servos and filament; the viewpoint model is naive, so these numbers are a floor for what one moving camera can do. \todo{final numbers from the grid}
\end{abstract}
